% !TeX root = review-main.tex

\begin{center}
{\LARGE\textbf{第1回 復習テスト}}
\end{center}

\vspace{0.3em}
\begin{flushright}
氏名：\underline{\hspace{15em}}
\end{flushright}

\vspace{0.4em}
\noindent\textbf{1. リスニング対策編}(4 $\times$ 10点)

\noindent\textbf{【問1】}次の慣用表現を英語に直しなさい。

\noindent「(道案内で)札幌駅までの行き方を教えていただけますか?」

\testblank[Could you tell me how to get to Sapporo Station?]{0.92\linewidth}
\vfill

\noindent\textbf{【問2】}次の慣用表現を英語に直しなさい。

\noindent「(お店で)いらっしゃいませ」

\testblank[May I help you?]{0.92\linewidth}
\vfill

\noindent\textbf{【問3】}次の英語を日本語に直しなさい。

\noindent「May I speak to Mary?」

\testblank[メアリーに代わっていただけますか?]{0.92\linewidth}
\vfill

\noindent\textbf{【問4】}次の英語を日本語に直しなさい。

\noindent「What would you like to order?」

\testblank[何をご注文でしょうか?]{0.92\linewidth}
\vfill

\noindent\textbf{2. 英作文対策編}(6 $\times$ 10点)

\noindent\textbf{【問5】}次の日本語を、主語と動詞を含む英文1文で表現しなさい。

\noindent「彼女は、私を幸せにしてくれます。」

\testblank[She makes me happy.]{0.92\linewidth}
\vfill

\noindent\textbf{【問6】}次の日本語を、主語と動詞を含む英文1文で表現しなさい。

\noindent「最後の試合に負けて、とても悲しかった。」

\testblank[I was very sad because (when , that) we lost the last game.]{0.92\linewidth}
\vfill

\noindent\textbf{【問7】}次の英作文の際によく使用される表現を英語に直しなさい。

\noindent「言い換えると」

\testblank[In other words, ]{0.92\linewidth}
\vfill

\noindent\textbf{【問8】}次の英作文の際によく使用される表現を英語に直しなさい。

\noindent「そのため」

\testblank[Therefore, as a result, ]{0.92\linewidth}
\vfill

\noindent\textbf{【問9】}次の日本語を、英作文しやすい表現（言い換え）に直しなさい。

\noindent「異文化を尊重することは大切だ。」

\testblank[他の国の文化について知ることは大切だ。]{0.92\linewidth}
\vfill

\noindent\textbf{【問10】}問9の表現を参考に、次の日本語を、主語と動詞を含む英文1文で表現しなさい。

\noindent「異文化を尊重することは大切だ。」

\testblank[It is important to learn about other cultures.]{0.92\linewidth}

\vfill
\begin{flushright}
\textbf{得点} \underline{\hspace{8em}} ／ 100点
\end{flushright}
